\documentclass[border=4pt]{standalone}
\usepackage{amsmath,amssymb,mathtools,xcolor,varwidth}
\begin{document}
\begin{varwidth}{28em}
\large\bfseries
Finally Newton circumscribes the parallelograms $aKbl$, $bLcm$, $cMdn$,
each rising above the curve to its upper vertex.
The difference between the circumscribed figure and the inscribed figure
is the sum of the little tops $Kl$, $Lm$, $Mn$, $Do$; since all bases are equal,
this sum is the single rectangle $ABla$, of breadth $AB$ and height $Aa$.
But as $AB$ is diminished $in\ infinitum$, this rectangle becomes less
than any given space, so by Lemma~I the inscribed and circumscribed figures
become ultimately equal, and $much\ more$ is the intermediate curvilinear figure
equal to either.
$Q.E.D.$
\end{varwidth}
\end{document}
